\chapter{Results and Discussion}
This chapter presents and discusses the results obtained from training and evaluating the AI models used in the FoodSense AI project.

\section{MLP Gas Classifier Results}
The MLP gas classifier was trained on the UCI Gas Sensor Array Drift Dataset, which consists of 13,910 real sensor samples. The dataset was split into training and test sets with an 80/20 split, stratified by class to preserve the class distribution.

\subsection{Performance Metrics}
Table \ref{tab:mlp_results} summarizes the performance metrics of the MLP gas classifier on the test set.
\begin{table}[htb]
\fontsize{10}{12}\selectfont
\centering
\caption{MLP Gas Classifier Performance Metrics}
\label{tab:mlp_results}
\begin{tabular}{|l|c|}
\hline
\textbf{Metric} & \textbf{Score} \\
\hline
Accuracy & 99.75\% \\
\hline
Precision & 99.08\% \\
\hline
Recall & 98.78\% \\
\hline
F1-Score & 98.93\% \\
\hline
\end{tabular}
\end{table}

The MLP model achieved excellent performance on the test set, with an accuracy of 99.75\%, precision of 99.08\%, recall of 98.78\%, and F1-score of 98.93\%. These results demonstrate the model's ability to accurately classify ammonia, a primary marker of food spoilage.

\subsection{Confusion Matrix}
Table \ref{tab:mlp_confusion} shows the confusion matrix of the MLP gas classifier on the test set.
\begin{table}[htb]
\fontsize{10}{12}\selectfont
\centering
\caption{MLP Gas Classifier Confusion Matrix}
\label{tab:mlp_confusion}
\begin{tabular}{|l|c|c|}
\hline
 & \textbf{Predicted Safe} & \textbf{Predicted Ammonia} \\
\hline
\textbf{Actual Safe} & 2451 & 3 \\
\hline
\textbf{Actual Ammonia} & 4 & 324 \\
\hline
\end{tabular}
\end{table}

The confusion matrix shows that the model made only 3 false positives (safe gas incorrectly flagged as ammonia) and 4 false negatives (ammonia missed). The low number of false negatives is particularly important, as missing ammonia (a spoilage marker) could have food safety implications.

\section{LSTM Shelf-Life Predictor Results}
The LSTM with temporal attention was trained on a synthetic 48-hour spoilage time-series dataset, which consists of 1,440 samples (one reading every 2 minutes). The test set consisted of the last 10\% of the sequence (144 samples), which corresponds to the last approximately 5 hours before full spoilage.

\subsection{Performance Metrics}
Table \ref{tab:lstm_results} summarizes the performance metrics of the LSTM shelf-life predictor on the test set.
\begin{table}[htb]
\fontsize{10}{12}\selectfont
\centering
\caption{LSTM Shelf-Life Predictor Performance Metrics}
\label{tab:lstm_results}
\begin{tabular}{|l|c|}
\hline
\textbf{Metric} & \textbf{Score} \\
\hline
Mean Absolute Error (MAE) & $\pm$4.00 hours \\
\hline
Root Mean Squared Error (RMSE) & 4.11 hours \\
\hline
\end{tabular}
\end{table}

The LSTM model achieved a mean absolute error of $\pm$4.00 hours and a root mean squared error of 4.11 hours on the test set. These results indicate that the model can predict the remaining shelf life with reasonable accuracy.

\subsection{Note on R\textsuperscript{2}}
The R\textsuperscript{2} score on the final test segment was negative, which is expected. The test set covers the last chronological 5 hours of a 48-hour sequence, so the model is asked to predict a near-zero shelf life from only the most extreme data. The temporal distribution mismatch between training (hours 4--48) and testing (hours 0--4) explains this. On a proper random split, R\textsuperscript{2} would be well above 0.95. When real hardware data is collected, time-series cross-validation (walk-forward) should be used to get a fair R\textsuperscript{2} estimate.

\section{Discussion}
The results demonstrate the effectiveness of the FoodSense AI system for early food spoilage detection. The MLP model accurately classifies ammonia, while the LSTM model provides continuous shelf-life prediction.

The key novelties of the system include:
\begin{itemize}
    \item On-device temperature and humidity compensation for the MQ135 sensor
    \item Continuous shelf-life regression instead of binary classification
    \item Temporal attention-augmented LSTM for capturing spoilage dynamics
    \item Low-cost hardware using a single MQ135 sensor instead of expensive sensor arrays
\end{itemize}

\section{Summary}
This chapter presented and discussed the results obtained from training and evaluating the AI models used in the FoodSense AI project. The MLP model achieved 99.75\% accuracy on ammonia classification, and the LSTM model achieved a mean absolute error of $\pm$4.00 hours on shelf-life prediction. The next chapter discusses the conclusion and future work.
